% !TEX root = ../enac_project_fr.tex
%%%%%%%%%%%%%%%%%%%%%%%%%%%%%%%%%%%%%%%%%%%%%%%%%%%%%%%%%%%%%%%%%%%%
% Chapter 5 - Implémentation en Java avec Orekit
%%%%%%%%%%%%%%%%%%%%%%%%%%%%%%%%%%%%%%%%%%%%%%%%%%%%%%%%%%%%%%%%%%%%

\chapter{Implémentation en Java avec Orekit}
\label{implementation_orekit}

Les chapitres précédents ont établi les modèles d'observation, les méthodes
d'estimation et les corrections nécessaires au PPP. Le présent chapitre décrit
leur traduction logicielle. Il
précise comment les fonctions disponibles dans Orekit ont été assemblées,
quelles adaptations ont été nécessaires et pour quelles raisons.

L'application est écrite en Java~17 et utilise Orekit~13.1.6~\cite{orekit131api}. Hipparchus~4.0.3~\cite{hipparchus403}
fournit les outils numériques employés indirectement par Orekit, notamment
l'optimisation non linéaire et le filtre de Kalman. Le principe directeur du
développement a été de conserver les objets et les algorithmes d'Orekit dès
qu'ils répondaient au besoin. Le code est ainsi limité aux règles
de traitement GNSS absentes de la bibliothèque.



\section{État de l'art des fonctions disponibles dans Orekit}
\label{sec:orekit_state_of_art}

Orekit est principalement orienté vers la mécanique spatiale et la
détermination d'orbite. Il possède néanmoins une infrastructure GNSS étendue :
lecture de produits, propagation des satellites, combinaisons d'observables,
modèles de mesure et estimateurs. La disponibilité d'une brique ne signifie
pas toujours qu'elle puisse être employée directement. Plusieurs modèles
physiques existent dans la bibliothèque sans disposer d'un
\texttt{EstimationModifier} compatible avec les mesures
\texttt{OneWayGNSSRange} et \texttt{OneWayGNSSPhase} retenues ici~\cite{orekit131api}.

Les tableaux~\ref{tab:orekit_data_support} et
\ref{tab:orekit_model_support} distinguent trois niveaux de prise en charge :
\emph{utilisation directe} lorsque l'API répond au besoin sans ajout
fonctionnel, \emph{adaptation nécessaire} lorsqu'un modèle ou un parseur
Orekit doit être raccordé aux conventions du traitement, et
\emph{développement spécifique} lorsque la fonction requise est absente.

\begin{table}[H]
\centering
\caption{Disponibilité des fonctions de lecture et de préparation des données.}
\label{tab:orekit_data_support}
\small
\begin{tabularx}{\textwidth}{p{3.3cm} p{3.0cm} X}
\hline
\textbf{Fonction} & \textbf{Support Orekit} & \textbf{Choix réalisé} \\
\hline
RINEX observation et navigation & Utilisation directe &
Lecture par les parseurs Orekit, y compris la décompression Hatanaka. \\

Orbites diffusées & Adaptation nécessaire &
Propagation assurée par Orekit ; ajout d'une politique de sélection du message
sain. \\

SP3 et CLK & Adaptation nécessaire &
Lecture et interpolation assurées par Orekit ; interface commune avec les
produits diffusés. \\

SINEX station & Adaptation nécessaire &
Lecture par le parseur Orekit ; filtrage préalable de blocs non exploitables
dans certains produits. \\

SINEX-BIAS & Adaptation nécessaire &
Lecture par le parseur Orekit ; application des OSB aux observables avec leur
convention d'unité et de signe. \\

ANTEX & Adaptation nécessaire &
Lecture des calibrations par Orekit ; composition des corrections mono- et
bi-fréquences dans les modèles de mesure. \\

Combinaisons GNSS & Utilisation directe &
Combinaisons iono-free, geometry-free et Melbourne--Wübbena fournies
par Orekit. \\

Sélection des signaux & Développement spécifique &
Définition des priorités entre codes RINEX et des règles de validation. \\
\hline
\end{tabularx}
\end{table}

\begin{table}[H]
\centering
\caption{Disponibilité des modèles physiques et des fonctions d'estimation.}
\label{tab:orekit_model_support}
\small
\begin{tabularx}{\textwidth}{p{3.3cm} p{3.0cm} X}
\hline
\textbf{Fonction} & \textbf{Support Orekit} & \textbf{Choix réalisé} \\
\hline
Code et phase GNSS unidirectionnels & Utilisation directe &
Réutilisation de \texttt{OneWayGNSSRange} et
\texttt{OneWayGNSSPhase}. \\

Temps de vol et rotation terrestre & Utilisation directe &
Calcul interne au modèle de mesure dans les repères Orekit ; aucune correction
de Sagnac n'est recodée. \\

Horloge satellite et TGD & Adaptation nécessaire &
Modèles d'horloge Orekit réutilisés ; adaptateurs propres pour unifier NAV et
CLK et appliquer le TGD avec la convention de l'observable. \\

Relativité et effet de Shapiro & Utilisation directe &
Réutilisation des modificateurs GNSS unidirectionnels fournis par Orekit quand
le mode de traitement les requiert. \\

Ionosphère et troposphère & Adaptation nécessaire &
Réutilisation de Klobuchar, Saastamoinen, Niell et
\texttt{EstimatedModel} ; modificateurs propres pour les deux classes de
mesure retenues. \\

Gradients troposphériques & Développement spécifique &
Extension de \texttt{EstimatedModel} par un modèle exposant deux
\texttt{ParameterDriver} supplémentaires. \\

PCO et PCV & Adaptation nécessaire &
Réutilisation d'ANTEX et du calcul de centre de phase Orekit ; assemblage
propre pour le récepteur terrestre et la combinaison iono-free. \\

Marées solides & Adaptation nécessaire &
Déplacement calculé par \texttt{TidalDisplacement} ; projection sur la ligne
de visée réalisée par un modificateur propre. \\

Wind-up & Adaptation nécessaire &
Spécialisation de \texttt{AbstractWindUp} avec l'attitude du satellite et le
repère topocentrique du récepteur. \\

ISB & Développement spécifique &
Ajout d'un \texttt{ParameterDriver} commun aux mesures Galileo. \\

Sauts de cycle & Adaptation nécessaire &
Détecteur geometry-free et combinaison MW d'Orekit complétés par le LLI, le
test de saut MW et la politique de réinitialisation. \\

WLS & Utilisation directe &
Réutilisation de \texttt{BatchLSEstimator} et de l'optimiseur
Levenberg--Marquardt d'Hipparchus. \\

EKF & Adaptation nécessaire &
Réutilisation de \texttt{KalmanEstimator} au moyen d'un propagateur de
récepteur statique et d'une gestion propre des états GNSS variables. \\
\hline
\end{tabularx}
\end{table}

Cette analyse conduit à une architecture d'extension plutôt qu'à une chaîne
de calcul indépendante. Les classes propres ne remplacent pas les modèles
d'Orekit : elles les rendent compatibles avec le type de mesure, les unités et
la structure d'état du positionnement au sol.


\section{Architecture logicielle}
\label{sec:software_architecture}

\subsection{Découpage en responsabilités}

Le code est organisé en couches aux responsabilités limitées. Cette séparation
évite que les règles de sélection des observables, les modèles physiques et la
logique d'estimation soient dupliqués. Elle permet aussi de remplacer une
source de données ou un solveur sans modifier les autres couches.

\begin{table}[H]
\centering
\caption{Responsabilités des principales couches de l'application.}
\label{tab:software_packages}
\small
\begin{tabularx}{\textwidth}{p{3.2cm} X}
\hline
\textbf{Couche} & \textbf{Responsabilité} \\
\hline
Configuration & Définition de la session, initialisation d'Orekit et
centralisation du modèle stochastique. \\
Données & Lecture, regroupement par époque, sélection des observables et
application des biais. \\
Satellites & Accès uniforme aux éphémérides, aux horloges et aux attitudes. \\
Mesures & Construction des observations Orekit et assemblage des corrections
physiques. \\
Propagation & Représentation du récepteur terrestre dans l'interface attendue
par les estimateurs Orekit. \\
Estimation & Orchestration des résolutions WLS et EKF, des ambiguïtés et des
covariances. \\
\hline
\end{tabularx}
\end{table}

On utilise un fichier de configuration au format JSON pour décrire la session
et sélectionner les modes de traitement. Les dépendances suivent le sens de la
chaîne suivante :

\begin{center}
\small
Configuration
$\longrightarrow$ données sélectionnées
$\longrightarrow$ états des satellites
$\longrightarrow$ mesures corrigées
$\longrightarrow$ estimation
$\longrightarrow$ résultats.
\end{center}

Le point d'entrée reste ainsi limité au chargement de la configuration, au
choix du solveur et au lancement du traitement. La construction des modèles
est déléguée aux couches concernées, ce qui évite de concentrer les choix
fonctionnels dans une procédure principale.

\subsection{Configuration des traitements}

Le fichier JSON permet de configurer une session sans modifier ni recompiler
le programme. Les principaux choix peuvent être définis indépendamment :

\begin{itemize}
    \item les données traitées : station, jour et constellations ;
    \item les observables : code seul ou code et phase, ainsi que la règle de
    sélection des signaux ;
    \item les modèles physiques : correction ionosphérique, modèle
    troposphérique et éphémérides diffusées ou précises ;
    \item la méthode de résolution : WLS ou EKF ;
    \item le modèle stochastique : pondération des mesures, incertitudes
    initiales, bruits de processus et seuils de rejet.
\end{itemize}

Le code a été conçu de manière modulaire afin que ces options puissent être
combinées librement dès lors que la configuration reste physiquement
cohérente. Par exemple, une combinaison iono-free exige deux
fréquences, le traitement de la phase impose la gestion des ambiguïtés et le
mode précis nécessite les produits d'orbite, d'horloge et de biais associés à
la session. Chaque option sélectionne une implémentation derrière une
interface commune.


\section{Lecture et préparation des données}
\label{sec:data_implementation}

\subsection{Initialisation d'Orekit et parsage}

Les parseurs Orekit lisent ensuite les observations, la navigation, la
position de référence et les calibrations d'antenne. En mode précis, ils
chargent également les orbites, les horloges et les biais observables. Les
produits sont résolus à partir du jour de la session afin de garantir leur
cohérence temporelle. Le filtre Hatanaka fourni par Orekit permet enfin de lire les
observations compactes sans étape de décompression externe.



\section{Construction des mesures Orekit}
\label{sec:measurement_implementation}

\subsection{Objets de mesure et combinaisons}

La couche de construction transforme chaque satellite retenu en mesure
Orekit. Les deux classes introduites au début de ce chapitre reçoivent
respectivement le code en mètres et la phase en cycles. Elles ont été conçues
pour une liaison unidirectionnelle entre un satellite GNSS émetteur et un
satellite récepteur dont l'état est estimé, par exemple dans une application
de détermination d'orbite. Elles conviennent néanmoins au positionnement au
sol si le récepteur est fourni sous la forme d'un état propagé compatible.

Ce choix permet de réaliser les fonctions suivantes : :

\begin{itemize}
    \item le calcul automatique des dérivées par rapport à l'état cartésien et
    au biais d'horloge du récepteur ;
    \item l'intégration des \texttt{ParameterDriver} de mesure ;
    \item la gestion d'une ambiguïté par l'\texttt{AmbiguityDriver} d'Orekit pour
    la phase.
\end{itemize}

En mode iono-free, les deux codes et les deux phases sont combinés par
\texttt{IonosphereFreeCombination}. La fréquence combinée fournie par Orekit
détermine la longueur d'onde synthétique utilisée par
\texttt{OneWayGNSSPhase}~\cite{orekit131api}.

\subsection{Assemblage par modificateurs}

Une mesure Orekit fournit le terme géométrique et les horloges associées, puis
accepte une liste de \texttt{EstimationModifier}. Chaque modificateur reçoit
la valeur théorique, les participants et, si nécessaire, les dérivées. Cette
interface permet d'ajouter une contribution sans modifier les classes de
mesure ni les solveurs.

Une fabrique unique centralise cet assemblage. Selon la configuration, elle
ajoute dans un ordre déterminé l'horloge satellite, la troposphère, les
marées, les termes relativistes, le wind-up, les antennes, l'ionosphère, l'ISB
et les filtres de résidus. Ce point d'entrée unique garantit qu'une correction
n'est ni oubliée ni appliquée deux fois selon le solveur.

\section{Extensions développées autour d'Orekit}
\label{sec:orekit_extensions}

\subsection{Horloges satellite, TGD et multi-constellation}

Les produits NAV et CLK n'exposent pas leur horloge sous la même forme. Une
interface commune les ramène au décalage d'horloge, au retard de groupe et à
l'indication de présence du terme relativiste. Les modificateurs de code et de
phase évaluent ensuite cette horloge à l'instant d'émission fourni par la
mesure.

Pour le traitement GPS--Galileo, un modificateur porte un unique
\texttt{ParameterDriver} d'ISB exprimé en mètres. Ses variantes code et phase
appliquent respectivement ce paramètre en mètres et en cycles, puis fournissent
au moteur d'estimation la dérivée correspondante. Toutes les mesures Galileo
partagent ce même driver ; les mesures GPS définissent l'horloge de référence.

\subsection{Modèles ionosphériques et troposphériques}

Les calculs physiques sont confiés aux modèles Orekit. Le développement propre
porte sur leur raccordement aux mesures unidirectionnelles :

\begin{itemize}
    \item le modèle de Klobuchar est appelé avec les signes et les unités
    propres au code et à la phase ;
    \item les modificateurs troposphériques propagent les dérivées des modèles
    Orekit vers les \texttt{ParameterDriver} estimés ;
    \item les modes déterministe et ZTD assemblent les modèles de Saastamoinen,
    la fonction de projection de Niell et le modèle estimé disponibles dans
    Orekit ;
    \item une extension du modèle ZTD ajoute les drivers des gradients nord et
    est. Seul le terme azimutal absent d'Orekit est développé ; les autres
    contributions restent calculées par les modèles d'Orekit.
\end{itemize}


\subsection{Corrections d'antenne}

Le lecteur ANTEX d'Orekit fournit les PCO, les PCV, leur période de validité et
le type d'antenne satellite. La bibliothèque contient également une base de
modificateur de centre de phase, mais pas l'assemblage complet requis pour une
station au sol et une observable iono-free.

Les modificateurs développés réalisent donc trois opérations :

\begin{enumerate}
    \item calculer la contribution satellite sur chaque fréquence avec la
    géométrie et l'attitude Orekit ;
    \item transformer le PCO récepteur du repère ANTEX nord--est--haut vers le
    repère topocentrique est--nord--haut, puis évaluer la PCV dans la direction
    du satellite ;
    \item combiner les deux corrections avec les coefficients de la
    combinaison iono-free et convertir le résultat en cycles pour la
    phase.
\end{enumerate}

\subsection{Marées solides et wind-up}

Le déplacement de marée est calculé par le modèle Orekit conforme aux
conventions IERS~2010.

De même, le traitement du wind-up spécialise la classe abstraite fournie par
Orekit. L'orientation du dipôle émetteur provient de l'attitude du satellite ;
celle du récepteur est définie par le repère topocentrique de la station. Un
état est conservé pour chaque satellite afin d'assurer la continuité entière
entre les époques. Pour une phase iono-free, un facteur d'échelle relie
le cycle physique à la longueur d'onde synthétique de la mesure.


\subsection{Détection des sauts de cycle}

La détection compose plusieurs briques au lieu de dupliquer les algorithmes
disponibles. La combinaison geometry-free et son ajustement polynomial sont
confiés au détecteur Orekit, tandis que la combinaison Melbourne--Wübbena est
calculée par la fabrique d'Orekit. Le code propre se limite à lire les bits LLI,
à comparer la variation MW au seuil et à fusionner les causes de détection par
satellite.

La décision de réinitialisation est séparée de la détection. Le gestionnaire
d'ambiguïtés applique la politique décrite à la
section~\ref{sec:ekf_gnss_features} : nouvelle valeur phase moins code lors de
la première apparition, d'un saut ou d'une interruption longue, puis
réinitialisation de la variance correspondante.


\section{Représentation des paramètres estimés}
\label{sec:estimated_parameters_implementation}

Orekit représente chaque inconnue ajustable par un \texttt{ParameterDriver}.
Celui-ci contient une valeur physique, une échelle de normalisation, une date
de référence et un indicateur de sélection. L'application emploie ce mécanisme
pour l'horloge récepteur, l'ISB Galileo, les ambiguïtés, le ZTD et les deux
gradients troposphériques.

Les ambiguïtés utilisent plus précisément \texttt{AmbiguityCache} et
\texttt{AmbiguityDriver}. Une clé constituée du satellite, du récepteur et de
la longueur d'onde garantit qu'une même ambiguïté est partagée par toutes les
mesures d'un arc. Leur valeur reste exprimée en cycles dans le driver, tandis
que les covariances configurées en mètres sont converties avec la longueur
d'onde propre à chaque signal.

Les matrices initiales et les bruits de processus sont fournis au filtre par
des implémentations de \texttt{CovarianceMatrixProvider}. Le bloc cartésien et
le bloc des paramètres de mesure sont construits séparément puis assemblés en
matrice bloc-diagonale. Les variances de marche aléatoire sont multipliées par
l'intervalle temporel entre deux époques. Cette organisation maintient la
cohérence entre les unités physiques de la configuration et les unités des
drivers Orekit.


\section{Résolution avec les estimateurs Orekit}
\label{sec:orekit_solvers}

\subsection{Moindres carrés pondérés}

Un estimateur par lots Orekit indépendant est construit pour chaque époque. La
position cartésienne est exposée par la représentation statique du récepteur,
tandis que l'horloge, l'ISB, les ambiguïtés et les paramètres troposphériques
sont apportés par leurs drivers. L'estimateur construit les résidus et les
jacobiennes, puis utilise l'optimiseur de Levenberg--Marquardt d'Hipparchus~\cite{hipparchus403}.

Le solveur ne forme donc pas explicitement la matrice normale de
l'équation~\eqref{eq:wls_solution}. Cette matrice résulte des dérivées fournies
par les mesures et leurs modificateurs. La tolérance, le nombre maximal
d'itérations et le nombre maximal d'évaluations restent configurables.

Une estimation indépendante par époque peut rendre le ZTD ou les gradients
peu observables. Le solveur contrôle alors les écarts-types obtenus. Si les
gradients ne sont pas déterminables, l'époque est recalculée avec le ZTD seul ;
si ce dernier reste lui-même inobservable, le modèle déterministe est utilisé.
Cette dégradation contrôlée évite de conserver une solution numérique ayant
convergé le long d'une direction presque singulière.

\subsection{Adaptation du filtre de Kalman d'Orekit}

\subsubsection{Adaptation d'un estimateur de détermination d'orbite}

\texttt{KalmanEstimator} est conçu pour la détermination d'orbite~\cite{orekit131api}. Il attend
un \texttt{PropagatorBuilder}, une trajectoire de référence et une matrice de
transition associée. Réimplémenter un EKF indépendant aurait dupliqué le
mécanisme de normalisation, le calcul des jacobiennes, le gain et la mise à
jour de covariance déjà fournis par Orekit et Hipparchus. Le choix a donc été
d'adapter la notion de propagation au récepteur fixe.

\texttt{StaticReceiverPropagator} conserve une position constante dans
l'ITRF. À chaque date, il la transforme dans le GCRF et l'encapsule dans un
\texttt{SpacecraftState}, type attendu par les modèles de mesure. L'orbite
cartésienne associée n'est pas interprétée comme une orbite physique ; elle
sert de conteneur compatible avec l'API d'estimation.

Cette couche d'adaptation remplit trois responsabilités : exposer les
coordonnées cartésiennes à l'estimateur, maintenir une position fixe en repère
terrestre entre deux corrections et fournir la matrice de transition. Cette
dernière transporte les corrections cartésiennes entre les bases GCRF des deux
époques par l'intermédiaire de l'ITRF.

Cette adaptation permet au moteur Orekit de conserver la position dans la
partie propagée de l'état, tandis que l'horloge, l'ISB, les ambiguïtés et la
troposphère restent des paramètres de mesure.

\subsubsection{Traitement d'une époque}

Le solveur orchestre chaque époque selon l'ordre suivant :

\begin{enumerate}
    \item construire les satellites valides et contrôler le GDOP ;
    \item détecter les sauts de cycle et préparer les ambiguïtés ;
    \item construire les phases puis les codes avec leurs modificateurs ;
    \item réestimer par code l'horloge et l'ISB lorsque leur réinitialisation
    est prévue ;
    \item reconstruire l'estimateur si la structure de l'état a changé ;
    \item évaluer les filtres de résidus avec l'état prédit ;
    \item regrouper les mesures acceptées dans une
    \texttt{MultiplexedMeasurement} et exécuter une mise à jour EKF unique ;
    \item enregistrer l'état, sa covariance et les ambiguïtés.
\end{enumerate}

Le regroupement des observations d'une même date évite de propager le
récepteur entre deux mesures simultanées. Le calcul du gain, de l'innovation
et de la covariance corrigée reste entièrement réalisé par
\texttt{KalmanEstimator}, qui s'appuie sur
\texttt{ExtendedKalmanFilter} d'Hipparchus~\cite{hipparchus403}.


\section{Sorties}
\label{sec:implementation_outputs}

Les deux solveurs produisent un fichier CSV de structure commune contenant la
date, le nombre de satellites, la position ECEF et géodétique, l'horloge, les
erreurs ENU, les écarts-types, l'ISB et les paramètres troposphériques. Lorsque
la phase est utilisée, un second fichier enregistre les ambiguïtés par
satellite. Le chemin de sortie encode le jour, la station, le mode
atmosphérique, les constellations, le solveur et le type d'éphémérides. Cette
convention rend chaque résultat traçable jusqu'à sa configuration.


Au total, l'implémentation repose sur Orekit pour les opérations génériques et
sensibles numériquement : temps, repères, propagation GNSS, interpolation des
produits précis, combinaisons, dérivées et estimation. Le code propre porte sur
la sémantique du traitement au sol : choix des observables, assemblage des
corrections, paramètres multi-constellation, continuité des ambiguïtés et
adaptation du filtre de détermination d'orbite. Cette frontière limite la
duplication de modèles validés tout en laissant explicites les choix propres au
positionnement PPP développé.
